% !TeX root = ../papers/nc-p3.tex
\section{位数4の半周期成分の二倍率と二つの二次式}
\label{sec:half-period-two-quadrics}
\begin{theorem}\label{thm:half-period-two-quadrics}
定理\ref{thm:half-period-components}の五次元成分は、有限な表示の選択を許せば、
二つの倍率上の $(2,2)$ 完全交叉三次元多様体の族で稠密開部分が覆われる。
倍率は $h^2-k^2=j^2-l^2=1$ とし $t=h^2-j^2\ne0$, $L=j^2+k^2$ とする。
四次元空間の基底で、以下の順に対角値を取る。
\[
 (d_i)=t^{-1}(jh,jh,-lk,-lk),\qquad
 (\lambda_i)=(jh,-jh,lk,-lk).
\]
平面 $U$ のPlücker座標 $p_{im}$（$1\le i<m\le4$）に対する方程式は
\[
 p_{12}p_{34}-p_{13}p_{24}+p_{14}p_{23}=0,\qquad
 \sum_{i<m}d_i d_m(\lambda_i+\lambda_m-L)p_{im}^2=0 .
\]
一般倍率で交叉は滑らかである。$\sum d_i d_m p_{im}^2\ne0$ 等の以下の開条件で、
平面の直交射影から元の四次テンソルを復元できる。
\end{theorem}
\begin{proof}[二つの交代部分の規格化]
前節の位数4の持上げを選び、$E_\pm,F_\pm$ の各二空間の基底を取る。
$E_+\otimes F_-$ と $E_-\otimes F_+$ 上の交代鉛筆の二つの階数1直線を
それぞれ $00,11$ に移す。第一切断の交代対角ブロックは
$00\wedge11$ の非零倍 $\alpha,\beta$ となる。
交代非対角ブロックは二平面間の $2\times2$ 行列であり、全交代部分が階数2である条件は
Pfaffianの消滅、すなわちその行列式が $\alpha\beta$ となることである。
第二切断 $N_{kj,il}=p_i M_{ij,kl}$ の二鉛筆は、第一切断の対応するブロックの
部分転置に符号を掛けたものなので、その像は $01,10$ の平面である。

四添字 $a=(a_1,a_2)\in\{0,1\}^2$ に $C a=(1-a_1,1-a_2)$ と $Ba=(a_2,a_1)$ を作用させ、
同じ記号を置換行列にも用いる。第一・第二の対称な八次ブロックの非対角 $4\times4$ 行列を
$H,J$ とする。鉛筆の台により $H$ の行 $a$ の非零位置は $CBa$、$J$ のそれは $Ba$ である。
非零値を $h_a,j_a$ と書く。第一の交代非対角行列は $(j_{uv})$、
第二は $(-h_{u,1-v})$ となる。第二の交代対角ブロックの係数は $-\alpha,\beta$ なので
\[
 h_{00}h_{11}-h_{01}h_{10}=j_{00}j_{11}-j_{01}j_{10}=\alpha\beta .
\]
逆に二つの交代部分はこのデータで全て復元できる。

$E_+,F_+,E_-,F_-$ の基底倍率を順に $x_i,y_i,u_i,v_i$ とすると、
$h_{ab}$ には $x_a v_{1-a}y_bu_{1-b}$、$j_{ab}$ には $x_a v_a y_bu_b$ が掛かる。
従って $\alpha=\beta=1$ の後でも二つの二行行列をそれぞれ独立に
行列式1の対角行・列倍率で規格化できる。実際 $x_0v_0,x_0v_1,y_0u_0,y_0u_1$ は
自由に選べ、相手添字の値は各々の逆となる。
全成分非零で $h_a=(h,k,k,h)$, $j_a=(j,l,l,j)$ にできる。
残る連続倍率は、作用しない倍率を除き二つの対称対角ブロックを互いに逆に拡大するものだけである。

これらの第一対角ブロックを $S,T$、$\Gamma$ を $2\times2$ 因子の第一位置での部分転置と書く。
$K=\operatorname{diag}(1,-1,-1,1)C$ とし、添字を $0,1,2,3$ の順にすれば
$G^\Gamma=G+\delta(G)K$、$\delta(G)=G_{12}-G_{03}$ が全対称行列で成立する。
両対称部分の階数4条件は
\[
 T=H^{\mathsf t}S^{-1}H,\qquad
 T^\Gamma=-J^{\mathsf t}(S^\Gamma)^{-1}J .
\]
両逆行列と $\delta(S)$ が非零の開部分を扱い、最後の倍率で $\delta(S)=1$ とする。
\end{proof}

\begin{proof}[対合の固有空間からの復元]
$\delta(T)=t=h^2-j^2=k^2-l^2$ と取る。この条件を課した族が五次元成分を支配することは
下で次元と交点から確認する。四次元空間とその双対の直和で、以下の写像を順に施す。
\[
 (x,y)\mapsto(H^{-1}y,H^{\mathsf t}x),\quad
 (x,y)\mapsto(x,y+tKx),\quad
 (x,y)\mapsto(J^{-\mathsf t}y,-Jx),\quad
 (x,y)\mapsto(x,y-Kx).
\]
合成 $F$ が $S$ のグラフを保てば二つのSchur条件が成立する。
$D=t^{-1}\operatorname{diag}(jh,-lk,-lk,jh)$、
$A=\operatorname{diag}(h/j,k/l,k/l,h/j)$、$R=-DAC$ と置く。
$F$ は反シンプレクティック対合で、その $\pm1$ 固有空間は $R\pm D$ のグラフである。
具体的に上側の二ブロックは $J^{-\mathsf t}H^{\mathsf t}$, $tJ^{-\mathsf t}KH^{-1}$、
下側はそれらへの左 $-K$ 倍と右下への $-JH^{-1}$ の和である。
$R\pm D$ のグラフに掛けて主張が従う。

正固有グラフに載せる水平平面を $U$ とする。負固有グラフに載せる平面を
$U^{\perp_D}$ とすれば両者の和はLagrange部分空間である。
$D|_U$ が非退化ならこれはグラフとなり、その対称行列は
\[
 S=R+D(2P_U-I),\qquad P_U=U(U^{\mathsf t}DU)^{-1}U^{\mathsf t}D .
\]
ここでは右辺の $U$ は平面の任意の列基底である。
$K$ によるトレースで $\delta(S)=1$ を書くと、
$\operatorname{tr}(D_{jh} C P_U)=L$ となる。ただし
$D_{jh}=\operatorname{diag}(jh,lk,lk,jh)$ である。
各相補添字対の和・差を $\sqrt2$ で割った基底で $D_{jh}C$ と $D$ は同時対角となり、
固有値は表記の $\lambda_i,d_i$ である。Cauchy--Binetの公式から
$\operatorname{tr}(D_{jh}CP_U)$ は
$\sum_{i<m}(\lambda_i+\lambda_m)d_i d_m p_{im}^2/\sum_{i<m}d_i d_m p_{im}^2$
であり、これが第二の二次式を与える。

$\delta(T)=t$ からは独立な方程式は増えない。実際、最初の写像は二つの固有グラフを
$t(R\mp D)$ のグラフに移し、正側の水平平面は $U'=H^{-1}(R+D)U$ となる。
よって $\delta(T)=t$ は $\operatorname{tr}(D_{jh}CP_{U'})=-L$ に等しい。
$U'$ の各固有座標の二乗の倍率は共通倍率を除いて
\[
 (z_i)=\big((j-h)^2,(j+h)^2,(l-k)^2,(l+k)^2\big)
\]
となる。置換 $CB$ は各固有線を保ち、そこでの追加符号は二乗で消える。
各 $i<m$ で
\[
 (\lambda_i+\lambda_m+L)z_i z_m=-t^2(\lambda_i+\lambda_m-L).
\]
同じ相補対なら $z_i z_m=t^2$、異なる対なら
$2L=(h^2+j^2)+(k^2+l^2)$ と $(j\pm h)(j\mp h)=(l\pm k)(l\mp k)=-t$ から従う。
従って両方のトレース条件はPlücker座標で同じ二次式となる。

二つの二次形式の鉛筆では、相補な二つのPlücker添字ごとに $2\times2$ に分かれる。
六つの判別根の平方は共通非零定数を除き
$L^2$, $L^2-(jh+lk)^2$, $L^2-(jh-lk)^2$ である。
一般にはこれらが相異なる非零値なので六根が単純であり、完全交叉は滑らかである。
滑らかな三次元完全交叉は連結ゆえ幾何学的に既約である。
\end{proof}

\begin{proof}[成分の被覆]
Walsh係数を行順に $a,b,c,d,e,f,g,v,w$ と書く。
持上げ $x=1$ で二鉛筆は既に対角線を持ち、規格化前の $H,J$ の重みは
$(c,-f,-f,c),(-v,g,-g,v)$ となる。$S$ の対角は $(b,d,d,b)$、
相補対の非対角は $a,e$、$T$ の対応値は $-b,d,d,-b$ と $a,-e$ となる。
従って $\delta(S)\delta(T)=a^2-e^2=c^2+v^2=f^2+g^2$ は符号平方の行・列和からの恒等式であり、
上で課した条件 $st=h^2-j^2$ はこのWalsh族上で成立する。
規格化前でも $F$ のグラフへの制限のトレースは零である。
上左ブロックの対角が零で、上右対角ブロックの相補二成分の符号が反対、
$S$ の該当対角値が等しいためである。従って固有次元は2ずつであり、
上の平面表示に入る。

53の共通点で $h^2,k^2,j^2,l^2$ は規格化後にそれぞれ $44,43,37,36$ であり、
各逆と六根の分離条件は全て単元である。
\path{computations/verify_half_period_two_quadrics.py} はこの点の固有次元と
一般倍率のグラフ・二次式の恒等式を検算する。
これで表示空間の一般ファイバーには必要な開条件を満たす点があり、総次元は5である。
二鉛筆の二根、残る四乗根倍率、持上げとその逆を除き基底は規格化されている。
$\delta(S)=1$ により有効に残る連続倍率もなく、$S$ から平面 $U$ は正固有グラフとの交叉で回収できる。
一般Walsh点の安定化群は有限なので、その近くでは持上げを忘れる繊維も有限である。
従ってこの既約表示の像は五次元で、一般Walsh点を通る位数4持上げの成分に属する。
定理\ref{thm:half-period-components}の一意性で、その大域的五次元成分を支配する。
\end{proof}
